\documentclass[11pt,a4paper]{article}
\usepackage[utf8]{inputenc}
\usepackage[T1]{fontenc}
% Langue : babel-french s'il est installé, sinon repli sans erreur
\IfFileExists{french.ldf}{%
  \usepackage[french]{babel}%
}{%
  \usepackage[english]{babel}%
  \addto\captionsenglish{\renewcommand{\contentsname}{Table des matières}}%
}
\usepackage{amsmath,amssymb,amsfonts}
\usepackage[margin=2.2cm]{geometry}
\usepackage{enumitem}
\usepackage{booktabs,array}
\usepackage[most]{tcolorbox}
\usepackage{hyperref}
\usepackage{fancyhdr}
\setlength{\parindent}{0pt}
\setlength{\parskip}{4pt}
\setlength{\headheight}{14pt}
\setlist{itemsep=1pt,topsep=2pt}
\pagestyle{fancy}\fancyhf{}
\fancyhead[L]{\small Chapitre 2 -- Fonctions génératrice et caractéristique}
\fancyhead[R]{\small 4ème DAS \& FinTech}
\fancyfoot[C]{\small\thepage}
\hypersetup{colorlinks=true,linkcolor=black}

\newcommand{\E}{\mathbb{E}}
\newcommand{\V}{\mathbb{V}}
\renewcommand{\P}{\mathbb{P}}
\newcommand{\N}{\mathbb{N}}
\newcommand{\R}{\mathbb{R}}
\newcommand{\C}{\mathbb{C}}
\newcommand{\loi}{\overset{\mathcal L}{=}}

\tcbset{defstyle/.style={colback=blue!4,colframe=blue!50!black,fonttitle=\bfseries,breakable,left=4pt,right=4pt,top=2pt,bottom=2pt}}
\newtcolorbox{defn}[1]{defstyle,title={Définition #1}}
\newtcolorbox{thm}[1]{defstyle,colback=red!3,colframe=red!50!black,title={#1}}
\newtcolorbox{rem}{defstyle,colback=gray!8,colframe=gray!60,title={Remarque}}

\begin{document}
\begin{center}
{\LARGE\bfseries Chapitre 2 : Fonctions génératrice et caractéristique\\[2pt] d'une variable aléatoire réelle}\\[6pt]
{\large Résumé de cours --- 4ème année DAS \& FINTECH (A.U. 2026-2027)}
\end{center}
\vspace{2pt}
\tableofcontents
\vspace{4pt}

\section{Introduction}
Les fonctions \textbf{génératrice} et \textbf{caractéristique} décrivent et étudient la loi d'une v.a.
\begin{itemize}
\item \textbf{Fonction génératrice} : proche d'un polynôme, ses coefficients sont les probabilités d'une v.a. \emph{discrète} à valeurs dans $\N$.
\item \textbf{Fonction caractéristique} : existe \emph{toujours} (discrète ou continue) ; elle résume la loi, donne les moments, et sert pour les sommes et la convergence vers la loi normale.
\end{itemize}
Applications : finance (risque, valorisation), assurance (sinistres, primes), IA (simulation et comparaison de distributions).

\section{Fonction génératrice}
\subsection{Définition et propriétés}
\begin{defn}{(fonction génératrice)}
Soit $X$ une v.a. à valeurs dans $\N$. Sa fonction génératrice est
\[ G_X(z)=\sum_{n\in X(\Omega)}\P(X=n)\,z^n=\E(z^X),\qquad -1\le z\le 1. \]
Elle ne dépend que de la loi de $X$. Si $X(\Omega)$ est fini, $G_X$ est un polynôme de degré $m=\max X(\Omega)$.
\end{defn}

\textbf{Propriétés.} $G_X(0)=\P(X=0)$ et $G_X(1)=1$. De plus $G_X$ est de classe $\mathcal C^\infty$ sur $[-1,1]$ et
\[ G_X^{(n)}(z)=\sum_{k=n}^{+\infty}\frac{k!}{(k-n)!}\,\P(X=k)\,z^{k-n}. \]

\begin{thm}{Proposition (caractérisation)}
$\displaystyle \P(X=n)=\frac{1}{n!}G_X^{(n)}(0)$ pour tout $n$ : la loi d'une v.a. discrète est entièrement caractérisée par $G_X$. Si $G_X^{(n)}(0)=0$ pour tout $n>N$, alors $X(\Omega)\subset\{0,\dots,N\}$.

\smallskip
\textbf{Théorème.} Pour $X,Y$ à valeurs dans $\N$ : $X\loi Y\iff G_X(z)=G_Y(z)\ \ \forall z\in[-1,1]$.
\end{thm}

\begin{thm}{Proposition (somme de v.a. indépendantes)}
Si $X$ et $Y$ sont indépendantes (à valeurs dans $\N$) : $G_{X+Y}(z)=G_X(z)\,G_Y(z)$.
La réciproque est \textbf{fausse}.
\end{thm}

\subsection{Fonctions génératrices des lois discrètes usuelles}
\begin{center}
\renewcommand{\arraystretch}{1.5}
\begin{tabular}{@{}llc@{}}
\toprule
Loi & $G_X(z)$ & Domaine\\
\midrule
Bernoulli $\mathcal B(p)$ & $(1-p)+pz$ & $[-1,1]$\\
Binomiale $\mathcal B(n,p)$ & $\big((1-p)+pz\big)^n$ & $[-1,1]$\\
Géométrique $\mathcal G(p)$ (rang du 1er succès) & $\dfrac{pz}{1-(1-p)z}$ & $|z|<\dfrac1{1-p}$\\
Poisson $\mathcal P(\lambda)$ & $e^{\lambda(z-1)}$ & $\R$\\
\bottomrule
\end{tabular}
\end{center}

\textbf{Démonstrations clés.}
\begin{itemize}
\item \emph{Binomiale} : $G_X(z)=\sum_{k=0}^n C_n^k(pz)^k(1-p)^{n-k}=((1-p)+pz)^n$ (binôme de Newton). On le retrouve aussi comme produit de $n$ fonctions génératrices de Bernoulli indépendantes.
\item \emph{Géométrique} : $\sum_{n\ge1}p(1-p)^{n-1}z^n=\frac{p}{1-p}\sum_{n\ge1}(z(1-p))^n=\frac{pz}{1-(1-p)z}$ (série géométrique).
\item \emph{Poisson} : $e^{-\lambda}\sum_{n\ge0}\frac{(\lambda z)^n}{n!}=e^{-\lambda}e^{\lambda z}$ (série exponentielle).
\end{itemize}

\subsection{Fonction génératrice et moments}
\begin{thm}{Théorème et corollaire}
$X$ (à valeurs dans $\N$) admet un moment d'ordre $p\in\N^*$ ssi $G_X$ est $p$ fois dérivable en $1$, et alors
\[ G_X^{(p)}(1)=\E\big[X(X-1)\cdots(X-p+1)\big]. \]
En particulier, si $X$ admet un moment d'ordre 2 :
\[ \E(X)=G_X'(1),\qquad \V(X)=G_X''(1)+G_X'(1)-\big(G_X'(1)\big)^2. \]
\end{thm}

\section{Fonction caractéristique}
\subsection{Définition et propriétés}
\begin{defn}{(fonction caractéristique)}
Pour une v.a. réelle $X$ :
\[ \Phi_X:\R\to\C,\qquad t\mapsto\Phi_X(t)=\E(e^{itX}). \]
\begin{itemize}
\item $X$ discrète : $\Phi_X(t)=\sum_{n\in X(\Omega)}e^{itn}\P(X=n)$ ;
\item $X$ continue de densité $f$ : $\Phi_X(t)=\int_{\R}e^{itx}f(x)\,dx$.
\end{itemize}
\end{defn}

\begin{thm}{Propriétés}
\begin{enumerate}[label=\arabic*.,nosep]
\item $\Phi_X(0)=1$ ;\quad 2. $|\Phi_X(t)|\le1$ pour tout $t$ ;
\item $\Phi_X$ est continue (convergence dominée) ;
\item $\Phi_X(-t)=\overline{\Phi_X(t)}$ ;
\item si $X$ est symétrique ($X\loi -X$), alors $\Phi_X$ est réelle (et paire).
\end{enumerate}
\end{thm}

\begin{thm}{Théorèmes et propositions}
\begin{itemize}
\item \textbf{Caractérisation} : $X\loi Y\iff\Phi_X(t)=\Phi_Y(t)\ \ \forall t\in\R$.
\item \textbf{Transformation affine} : $\Phi_{aX+b}(t)=e^{ibt}\,\Phi_X(at)$.
\item \textbf{Somme} : si $X\perp\!\!\!\perp Y$, alors $\Phi_{X+Y}(t)=\Phi_X(t)\,\Phi_Y(t)$.
\item \textbf{Inversion} : si $\int_\R|\Phi_X(t)|\,dt<+\infty$, alors $X$ admet la densité
$\displaystyle f(x)=\frac1{2\pi}\int_{\R}e^{-itx}\Phi_X(t)\,dt.$
\end{itemize}
\end{thm}

\subsection{Fonctions caractéristiques des lois usuelles}
\begin{center}
\renewcommand{\arraystretch}{1.6}
\begin{tabular}{@{}lc@{}}
\toprule
Loi & $\Phi_X(t)$\\
\midrule
Bernoulli $\mathcal B(p)$ & $(1-p)+pe^{it}$\\
Binomiale $\mathcal B(n,p)$ & $(1-p+pe^{it})^n$\\
Géométrique $\mathcal G(p)$ & $\dfrac{pe^{it}}{1-(1-p)e^{it}}$ \ (rang du 1er succès)\\
Poisson $\mathcal P(\lambda)$ & $e^{\lambda(e^{it}-1)}$\\
Uniforme $\mathcal U([0,1])$ & $\dfrac{e^{it}-1}{it}$ si $t\ne0$, $1$ si $t=0$\\
Uniforme $\mathcal U([a,b])$ & $\dfrac{e^{itb}-e^{ita}}{it(b-a)}$ si $t\ne0$, $1$ si $t=0$\\
Exponentielle $\mathcal E(1)$ & $\dfrac{1}{1-it}$\\
Exponentielle $\mathcal E(\lambda)$ & $\dfrac{\lambda}{\lambda-it}$\\
Cauchy $\mathcal C(0,1)$ & $e^{-|t|}$\\
Cauchy $\mathcal C(x_0,a)$ & $e^{ix_0t}e^{-a|t|}$\\
Normale $\mathcal N(0,1)$ & $e^{-t^2/2}$\\
Normale $\mathcal N(m,\sigma^2)$ & $e^{imt}\,e^{-\sigma^2t^2/2}$\\
\bottomrule
\end{tabular}
\end{center}
\emph{Note :} dans l'annexe du cours, la géométrique y est définie comme le nombre d'échecs avant le premier succès (sans le facteur $e^{it}$) ; ici on garde la convention du chapitre~1 ($\N^*$).

\textbf{Méthode.} On se ramène à la loi standard grâce à $\Phi_{aX+b}(t)=e^{ibt}\Phi_X(at)$ :
$\mathcal U([a,b])\loi(b-a)X+a$ avec $X\sim\mathcal U([0,1])$ ; $\mathcal E(\lambda)\loi\frac1\lambda X$ avec $X\sim\mathcal E(1)$ ; $\mathcal C(x_0,a)\loi aX+x_0$ ; $\mathcal N(m,\sigma^2)\loi\sigma Z+m$.

\begin{rem}
La fonction caractéristique de la loi de Cauchy n'est pas dérivable en $0$ : la loi de Cauchy n'admet pas d'espérance.
\end{rem}

\subsection{Fonction caractéristique et moments d'ordre supérieur}
\begin{thm}{Proposition et corollaire}
Si $X$ admet un moment d'ordre $p\in\N^*$, alors $\Phi_X$ est $p$ fois dérivable et
\[ \E(X^p)=\frac{1}{i^p}\,\Phi_X^{(p)}(0). \]
En particulier (moment d'ordre 2) : $\E(X)=-i\,\Phi_X'(0)$ et $\V(X)=-\Phi_X''(0)+\big(\Phi_X'(0)\big)^2$.

Si $\Phi_X$ n'est pas dérivable en $0$, alors $X$ n'admet pas de moment d'ordre 1.
\end{thm}

\section{Formulaire de synthèse}
\begin{center}
\renewcommand{\arraystretch}{1.4}
\begin{tabular}{@{}p{5.2cm}p{9cm}@{}}
\toprule
\textbf{Outil} & \textbf{À retenir}\\
\midrule
Fonction génératrice & $G_X(z)=\E(z^X)$, $X$ à valeurs dans $\N$ ; $G_X(0)=\P(X=0)$, $G_X(1)=1$\\
Lois $\leftrightarrow G_X$ & $\P(X=n)=G_X^{(n)}(0)/n!$ ; $X\loi Y\iff G_X=G_Y$\\
Moments via $G_X$ & $\E(X)=G_X'(1)$ ; $\V(X)=G_X''(1)+G_X'(1)-G_X'(1)^2$\\
Fonction caractéristique & $\Phi_X(t)=\E(e^{itX})$, existe toujours, $\Phi_X(0)=1$, $|\Phi_X|\le1$\\
Lois $\leftrightarrow\Phi_X$ & $X\loi Y\iff\Phi_X=\Phi_Y$ ; inversion si $\Phi_X\in L^1$\\
Moments via $\Phi_X$ & $\E(X^p)=\Phi_X^{(p)}(0)/i^p$ ; $\E(X)=-i\Phi_X'(0)$\\
Affine & $\Phi_{aX+b}(t)=e^{ibt}\Phi_X(at)$\\
Somme indépendante & $G_{X+Y}=G_XG_Y$ ; $\Phi_{X+Y}=\Phi_X\Phi_Y$ (réciproques fausses)\\
\bottomrule
\end{tabular}
\end{center}

\textbf{Formules utiles (annexe).} Série géométrique : $\sum_{n\ge0}x^n=\frac1{1-x}$ ($|x|<1$), $\sum_{n\ge1}nx^{n-1}=\frac1{(1-x)^2}$, $\sum_{n\ge2}n(n-1)x^{n-2}=\frac2{(1-x)^3}$. Série exponentielle : $e^x=\sum_{n\ge0}\frac{x^n}{n!}$. Combinaison : $C_n^k=\frac{n!}{k!(n-k)!}$. Binôme : $(a+b)^n=\sum_{k=0}^nC_n^ka^kb^{n-k}$.

\end{document}
